\nonstopmode{}
\documentclass[a4paper]{book}
\usepackage[times,inconsolata,hyper]{Rd}
\usepackage{makeidx}
\makeatletter\@ifl@t@r\fmtversion{2018/04/01}{}{\usepackage[utf8]{inputenc}}\makeatother
% \usepackage{graphicx} % @USE GRAPHICX@
\makeindex{}
\begin{document}
\chapter*{}
\begin{center}
{\textbf{\huge Package `thresholdanalysis'}}
\par\bigskip{\large \today}
\end{center}
\ifthenelse{\boolean{Rd@use@hyper}}{\hypersetup{pdftitle = {thresholdanalysis: Threshold Effects in Regression Models}}}{}
\ifthenelse{\boolean{Rd@use@hyper}}{\hypersetup{pdfauthor = {Song Sheng}}}{}
\begin{description}
\raggedright{}
\item[Title]\AsIs{Threshold Effects in Regression Models}
\item[Version]\AsIs{1.0.0}
\item[Description]\AsIs{
Provides a unified framework for detecting and modeling threshold effects 
in logistic regression, linear regression, and Cox proportional hazards models.
The package implements automated threshold identification using restricted cubic splines (RCS) 
and generalized additive models (GAM) with profile likelihood optimization.
Key features include:
- Confounder-adjusted analyses with 95% confidence intervals for threshold estimates
- Visualization of odds ratios (OR), hazard ratios (HR), and beta coefficients with 95% CI bands
- Excel-ready output of model coefficients and likelihood ratio test results
- Segmented regression modeling via X1/X2 variables for piecewise linear/nonlinear effects
Primary applications encompass dose-response relationship analysis, 
exposure threshold determination in clinical research, and epidemiological studies.}
\item[License]\AsIs{GPL-3}
\item[Encoding]\AsIs{UTF-8}
\item[Roxygen]\AsIs{list(markdown = TRUE)}
\item[Imports]\AsIs{dplyr, ggplot2, survival, openxlsx, broom, lmtest, mgcv,
purrr, readr, readxl, rms, scales, splines}
\item[Suggests]\AsIs{knitr, rmarkdown}
\item[Config/roxygen2/version]\AsIs{8.0.0}
\item[NeedsCompilation]\AsIs{no}
\item[Author]\AsIs{Song Sheng [aut, cre]}
\item[Maintainer]\AsIs{Song Sheng }\email{747367050@qq.com}\AsIs{}
\end{description}
\Rdcontents{Contents}
\HeaderA{cox\_threshold}{Cox Proportional Hazards Model with Threshold Effect Analysis}{cox.Rul.threshold}
%
\begin{Description}
This function performs a threshold effect analysis using Cox proportional hazards models,
including restricted cubic spline (RCS) modeling for the focal variable,
threshold determination, and model comparison \strong{adjusted for confounders}.
HR+95\%CI visualization with explicit threshold marking.
Outputs include Excel results and analysis objects.
\end{Description}
%
\begin{Usage}
\begin{verbatim}
cox_threshold(formula, data, var_name, rcs_nodes = 4)
\end{verbatim}
\end{Usage}
%
\begin{Arguments}
\begin{ldescription}
\item[\code{formula}] A formula specifying the model \strong{including both focal variable (var\_name) and confounders}.
Example: \code{Surv(time, status) \textasciitilde{} focal\_var + confounder1 + confounder2}

\item[\code{data}] Data frame containing all variables in the formula

\item[\code{var\_name}] Name of the focal variable to analyze for threshold effects

\item[\code{rcs\_nodes}] Number of knots for RCS (default=4) applied to focal variable
\end{ldescription}
\end{Arguments}
%
\begin{Value}
List containing:
\begin{itemize}

\item{} cutoff: Optimal threshold value
\item{} mdl0: Original model with all predictors (focal + confounders)
\item{} mdl1: Segmented model (X1/X2 for focal variable + confounders)
\item{} mdl2: Threshold test model (focal variable + X2 + confounders)
\item{} plot: Visualization object of HR curve with threshold marker

\end{itemize}

\end{Value}
\HeaderA{linear\_threshold}{Linear Regression with Threshold Effect Analysis}{linear.Rul.threshold}
%
\begin{Description}
This function performs a threshold effect analysis using linear regression models,
including restricted cubic spline (RCS) modeling for the focal variable,
threshold determination, and model comparison \strong{adjusted for confounders}.
Regression coefficient+95\%CI visualization with explicit threshold marking.
Outputs include Excel results and analysis objects.
\end{Description}
%
\begin{Usage}
\begin{verbatim}
linear_threshold(formula, data, var_name, rcs_nodes = 4)
\end{verbatim}
\end{Usage}
%
\begin{Arguments}
\begin{ldescription}
\item[\code{formula}] A formula specifying the model \strong{including both focal variable (var\_name) and confounders}.
Example: \code{outcome \textasciitilde{} focal\_var + confounder1 + confounder2}

\item[\code{data}] Data frame containing all variables in the formula

\item[\code{var\_name}] Name of the focal variable to analyze for threshold effects

\item[\code{rcs\_nodes}] Number of knots for RCS (default=4) applied to focal variable
\end{ldescription}
\end{Arguments}
%
\begin{Value}
List containing:
\begin{itemize}

\item{} cutoff: Optimal threshold value
\item{} mdl0: Original model with all predictors
\item{} mdl1: Segmented model (X1/X2 for focal variable + confounders)
\item{} mdl2: Threshold test model (focal variable + X2 + confounders)
\item{} plot: Visualization object of coefficient curve with threshold marker

\end{itemize}

\end{Value}
\HeaderA{logistic\_threshold}{Logistic Regression with Threshold Effect Analysis}{logistic.Rul.threshold}
%
\begin{Description}
This function performs a threshold effect analysis using logistic regression,
including restricted cubic spline (RCS) modeling for the focal variable,
threshold determination, and model comparison \strong{adjusted for confounders}.
OR+95\%CI visualization with explicit threshold marking.
Outputs include Excel results and analysis objects.
\end{Description}
%
\begin{Usage}
\begin{verbatim}
logistic_threshold(formula, data, var_name, rcs_nodes = 4)
\end{verbatim}
\end{Usage}
%
\begin{Arguments}
\begin{ldescription}
\item[\code{formula}] A formula specifying the model \strong{including both focal variable (var\_name) and confounders}.
Example: \code{outcome \textasciitilde{} focal\_var + confounder1 + confounder2}

\item[\code{data}] Data frame containing all variables in the formula

\item[\code{var\_name}] Name of the focal variable to analyze for threshold effects

\item[\code{rcs\_nodes}] Number of knots for RCS (default=4) applied to focal variable
\end{ldescription}
\end{Arguments}
%
\begin{Value}
List containing:
\begin{itemize}

\item{} cutoff: Optimal threshold value
\item{} mdl0: Original model with all predictors (focal + confounders)
\item{} mdl1: Segmented model (X1/X2 for focal variable + confounders)
\item{} mdl2: Threshold test model (focal variable + X2 + confounders)
\item{} plot: Visualization object of OR curve with threshold marker

\end{itemize}

\end{Value}
\printindex{}
\end{document}
